% Ch18 WS4 -- electrolysis, Faraday's law, and a 10-point FRQ. Block 5.
\documentclass{apchem-worksheet}

\apcunitword{Chapter}
\apcunit{18}
\apcunittitle{Electrochemistry}
\apcdoctitle{Worksheet 4 \textbullet\ Electrolysis, Faraday's Law \& FRQ}
\apcsubtitle{Zumdahl \S18.7--18.8 \textbullet\ the electron is a reagent; count it}

\begin{document}
\wsheader[$I = q/t$ \textbullet\ $F = \SI{96485}{\coulomb\per\mole}$
e$^-$ \textbullet\ time in \textbf{seconds}. Molar masses
(\si{\gram\per\mole}): \ce{Ag} 107.87 \textbullet\ \ce{Cu} 63.55
\textbullet\ \ce{Cr} 52.00 \textbullet\ \ce{Al} 26.98 \textbullet\
\ce{Ni} 58.69 \textbullet\ \ce{Au} 196.97.]

\problem[]{Write the conversion chain from a measured current to a mass of
plated metal, naming each step.
\answerlines[4]{Current and time give charge, $q = It$. Charge divided by
Faraday's constant gives moles of electrons. Moles of electrons divided by
the number of electrons in the half-reaction gives moles of metal. Moles
times molar mass gives the mass.}}

\problem[]{How many moles of electrons are needed to deposit one mole of
each metal?}
\begin{parts}
  \item \ce{Ag} from \ce{Ag+}: \answerblank[1.8cm]{1}
  \item \ce{Cu} from \ce{Cu^2+}: \answerblank[1.8cm]{2}
  \item \ce{Al} from \ce{Al^3+}: \answerblank[1.8cm]{3}
  \item \ce{Cr} from \ce{Cr^3+}: \answerblank[1.8cm]{3}
  \item Where does that number come from?
        \answerblank[9cm]{the charge on the ion, via the half-reaction}
\end{parts}

\problem[]{A current of \SI{2.00}{\ampere} passes through \ce{AgNO3}
solution for \SI{30.0}{\minute}.}
\begin{parts}
  \item Convert the time to seconds. \answerblank[2.4cm]{1800 s}
  \item Calculate the charge passed. \answerblank[2.4cm]{3600 C}
  \item Calculate the moles of electrons. \answerblank[2.6cm]{0.0373 mol}
  \item Calculate the mass of silver deposited. \workspace[2cm]
        \keyonly{\begin{apcanswer}\ce{Ag+ + e^- -> Ag}, so
        mol Ag $= 0.0373$;
        $m = (0.0373)(107.87) =
        \mathbf{4.03}~\text{g}$\end{apcanswer}}
\end{parts}

\problem[]{The same charge is now passed through \ce{CuSO4} solution
instead --- \SI{1.50}{\ampere} for \SI{1.00}{\hour}.}
\begin{parts}
  \item Calculate the mass of copper deposited. \workspace[2.6cm]
        \keyonly{\begin{apcanswer}$q = (1.50)(3600) = 5400$~C;
        mol e$^- = 5400/96485 = 0.0560$;
        \ce{Cu^2+ + 2e^- -> Cu} so mol Cu $= 0.0280$;
        $m = (0.0280)(63.55) =
        \mathbf{1.78}~\text{g}$\end{apcanswer}}
  \item A student forgets to divide by 2 and reports 3.56 g. By what factor
        is that wrong, and what step did they skip?
        \answerlines[2]{Wrong by a factor of 2. They converted moles of
        electrons directly into moles of copper, ignoring that the
        half-reaction requires two electrons per copper atom.}
\end{parts}

\problem[]{Working backwards.}
\begin{parts}
  \item How long, in hours, does it take to deposit \SI{5.00}{\gram} of
        copper at \SI{2.00}{\ampere}? \workspace[2.8cm]
        \keyonly{\begin{apcanswer}mol Cu $= 5.00/63.55 = 0.0787$;
        mol e$^- = 0.157$; $q = (0.157)(96485) = 15{,}183$~C;
        $t = 15183/2.00 = 7591$~s $= \mathbf{2.11}$~h\end{apcanswer}}
  \item What current is needed to deposit \SI{10.0}{\gram} of silver in
        exactly \SI{1.00}{\hour}? \workspace[2.6cm]
        \keyonly{\begin{apcanswer}mol Ag $= 10.0/107.87 = 0.0927$;
        mol e$^- = 0.0927$; $q = 8945$~C;
        $I = 8945/3600 = \mathbf{2.48}~\text{A}$\end{apcanswer}}
  \item Chromium is plated from \ce{Cr^3+} at \SI{5.00}{\ampere} for
        \SI{2.00}{\hour}. Find the mass. \workspace[2.6cm]
        \keyonly{\begin{apcanswer}$q = (5.00)(7200) = 36{,}000$~C;
        mol e$^- = 0.373$; mol Cr $= 0.373/3 = 0.124$;
        $m = (0.124)(52.00) = \mathbf{6.47}~\text{g}$\end{apcanswer}}
\end{parts}

\problem[]{Two electrolytic cells are wired in series, so the same current
passes through both for the same time. One contains \ce{AgNO3}, the other
\ce{CuSO4}.}
\begin{parts}
  \item Which receives more charge?
        \answerblank[6.4cm]{the same --- they are in series}
  \item Which deposits more \emph{moles} of metal, and why?
        \answerlines[3]{The silver cell, by a factor of two. Equal charge
        means equal moles of electrons delivered to each, but silver needs
        one electron per atom while copper needs two, so twice as many
        moles of silver are produced.}
  \item Does that mean the silver electrode gains more \emph{mass}?
        Explain. \answerlines[2]{Yes, and by more than a factor of two:
        silver has both the larger mole count and the larger molar mass
        (107.87 against 63.55).}
\end{parts}

\problem[]{Industrial scale.}
\begin{parts}
  \item How much charge is required to produce \SI{1.00}{\kilo\gram} of
        aluminium from \ce{Al^3+}? \workspace[2.6cm]
        \keyonly{\begin{apcanswer}mol Al $= 1000/26.98 = 37.1$;
        mol e$^- = 3(37.1) = 111$;
        $q = (111)(96485) =
        \mathbf{1.07\times10^{7}}~\text{C}$\end{apcanswer}}
  \item Explain, using two features of aluminium, why smelting it is so
        energy-intensive. \answerlines[3]{Each aluminium atom requires
        three electrons rather than one or two, and the \ce{Al^3+}/\ce{Al}
        reduction potential is very negative ($-1.66$~V), so a large
        applied voltage is needed as well as a large charge. Energy is
        charge times voltage, and both factors are large.}
  \item Why does recycling aluminium save so much energy compared with
        producing it? \answerlines[2]{Melting and reforming existing
        aluminium metal requires no electrolysis at all, so the enormous
        electrical energy of reducing \ce{Al^3+} is avoided entirely.}
\end{parts}

\problem[]{\textbf{FRQ (10 points).} An electrolytic cell is used to plate
nickel onto a steel part from a solution of \ce{NiSO4}. A current of
\SI{3.00}{\ampere} is passed for \SI{45.0}{\minute}.}
\begin{parts}
  \item Write the half-reaction occurring at the cathode.
        \answerblank[6cm]{\ce{Ni^2+ + 2e^- -> Ni}}
  \item Explain why the part being plated must be the cathode.
        \answerlines[3]{Plating means depositing nickel metal, which
        requires reducing \ce{Ni^2+} to \ce{Ni}. Reduction always occurs at
        the cathode, in electrolytic cells exactly as in galvanic ones, so
        the part must be connected as the cathode.}
  \item Calculate the mass of nickel deposited. \workspace[3cm]
        \keyonly{\begin{apcanswer}$t = 45.0 \times 60 = 2700$~s;
        $q = (3.00)(2700) = 8100$~C;
        mol e$^- = 8100/96485 = 0.0840$;
        mol Ni $= 0.0840/2 = 0.0420$;
        $m = (0.0420)(58.69) =
        \mathbf{2.46}~\text{g}$\end{apcanswer}}
  \item Given $E^\circ$ for \ce{Ni^2+/Ni} is $-0.23$~V and for the anode
        reaction $+1.23$~V, calculate $E^\circ_{\text{cell}}$ and state
        what it tells you about the cell. \workspace[2.4cm]
        \keyonly{\begin{apcanswer}$E^\circ_{\text{cell}} = -0.23 - 1.23 =
        \mathbf{-1.46}~\text{V}$. Negative, so the reaction is
        thermodynamically unfavored and the cell is
        \textbf{electrolytic} --- it requires an external applied potential
        of more than 1.46 V.\end{apcanswer}}
  \item Calculate $\Delta G^\circ$ for the cell reaction, taking $n = 2$.
        \workspace[2.2cm]
        \keyonly{\begin{apcanswer}$\Delta G^\circ = -(2)(96485)(-1.46) =
        +281{,}736~\text{J} = \mathbf{+282}~\text{kJ}$ --- positive, as
        expected for an unfavored process\end{apcanswer}}
\end{parts}

\begin{apcnote}[Rubric --- score yourself, 10 points]
\textbf{(a) 1 pt} balanced half-reaction with 2 electrons \textbullet\
\textbf{(b) 2 pts} 1 identifies that plating requires reduction, 1 states
reduction occurs at the cathode in \emph{all} cells --- an answer relying on
the electrode's sign earns 0 \textbullet\ \textbf{(c) 4 pts} 1 time
converted to seconds, 1 charge, 1 moles of electrons divided by 2, 1 mass
$= 2.46$~g \textbullet\ \textbf{(d) 2 pts} 1 $E^\circ = -1.46$~V, 1
identifies the cell as electrolytic and requiring an applied potential
\textbullet\ \textbf{(e) 1 pt} $\Delta G^\circ = +282$~kJ with the correct
positive sign.
\end{apcnote}

\begin{spiralstrip}{Chapter 18 \textbullet\ blocks 1--4}
  \item At equilibrium a cell has $E = $ \answerblank[1.6cm]{0}
  \item Standard conditions correspond to $Q = $ \answerblank[1.6cm]{1}
  \item In a concentration cell, the anode is the:
        \answerblank[3cm]{dilute side}
  \item $\Delta G^\circ = -nFE^\circ$; if $E^\circ < 0$ then $\Delta G^\circ$
        is \answerblank[2cm]{positive}
  \item Le Ch\^atelier applies to a running cell?
        \answerblank[1.6cm]{no}
\end{spiralstrip}

\end{document}
